\documentclass{article}
\usepackage{geometry}
\geometry{a4paper, margin=1in}
\title{ML Research Proposal: Inverse Design of Morphogenetic Patterns via Deep Surrogate Modeling}
\author{Candidate Name}
\date{Suggested Supervisors: Prof. Dr. Ivo F. Sbalzarini, Dr. Carl D. Modes}

\begin{document}
\maketitle

\section*{Abstract}
The concept of a "Compiler of Morphogenetic Stimuli" (Levin, 2012) proposes that a high-level biological "code" dictates anatomical outcome. Mechanistically, this code translates to physical inputs like active forces, cell polarity, and bioelectric gradients. However, predicting the necessary initial biophysical stimuli (the "code") required to generate a specific final tissue shape remains an intractable inverse problem due to the complexity of non-linear physics. This project aims to solve this challenge by developing a \textbf{Deep Learning Surrogate Model} for the \textbf{TopoSPAM} multiscale simulation platform \cite{1}. We will train a deep neural network (DNN) or Physics-Informed Neural Network (PINN) on a large dataset of TopoSPAM simulations \cite{1}, mapping high-dimensional input parameter space (e.g., spontaneous strain tensor, active polarity coefficients) to 3D shape outputs (e.g., mesh topology). The final model will function as a fast, searchable surrogate, enabling rapid \textbf{Inverse Design}: given a desired final tissue geometry (e.g., a perfect organoid), the ML model will instantly predict the biophysical parameters needed to initiate the TopoSPAM Active Vertex or Spring-Lattice models to produce that structure. This approach drastically accelerates the search for morphogenetic control programs.

\section*{Related Work}
\begin{itemize}
    \item \textbf{TopoSPAM:} Topology-grounded simulation of morphogenesis and active matter, providing the mechanistic ground truth data \cite{1}.
    \item \textbf{Levin's Compiler Concept:} The theoretical framework for reprogramming anatomy by defining an interface for morphogenetic instructions (Levin, 2012).
    \item \textbf{Physics-Informed Neural Networks (PINNs):} Use of ML to embed and accelerate complex physical simulation (e.g., active hydrodynamics) by encoding differential equations into the loss function (e.g., Raissi et al., 2019).
\end{itemize}

\end{document}
